% ============================================================
%  PLAN DE CONDUITE DU CHANGEMENT — Phase 2 : Planification
%  Time'Eats — Cellule PMO
% ============================================================
\documentclass[11pt,a4paper,oneside]{report}
\usepackage{../style/consulting}
\hypersetup{pdftitle={Plan de Conduite du Changement — Time'Eats PMO}}
\pagestyle{fancy}\fancyhf{}
\renewcommand{\headrulewidth}{0pt}
\fancyhead[L]{\begin{tikzpicture}[remember picture,overlay]
  \draw[cnNavy,line width=0.5pt](0,0)--(\textwidth,0);
\end{tikzpicture}\vspace{2pt}\timeeatslogo\hspace{4pt}\small\textcolor{cnSlate}{Time'Eats \textbf{·} Conduite du Changement — \textit{[Nom projet]}}}
\fancyhead[R]{\small\textcolor{cnSlate}{\thepage}}
\fancyfoot[C]{\tiny\textcolor{cnSlate}{Document confidentiel — Cellule PMO — CESI MAALSI 2026}}
\begin{document}\small

\begin{center}
{\LARGE\bfseries\color{cnNavy} Plan de Conduite du Changement}\\[4pt]
\textcolor{cnOrange}{\rule{10cm}{1pt}}\\[4pt]
{\large\color{cnSlate} Phase 2 — Planification \quad|\quad Projet : \textit{[Nom du projet]}}
\end{center}

\vspace{6pt}

\begin{decision}
La conduite du changement est un levier de succès critique. Un outil déployé non adopté = projet échoué. Ce plan garantit l'accompagnement des utilisateurs tout au long du projet.
\end{decision}

\section*{1 — Analyse de l'impact sur les utilisateurs}

\renewcommand{\arraystretch}{1.4}
\begin{tabularx}{\textwidth}{L{4cm} L{3cm} Y C{2.5cm}}
\toprule
\rowcolor{cnNavy}\th{Population} & \th{Effectif concerné} & \th{Impact du changement} & \th{Niveau résistance} \\
\midrule
\textit{[Équipe 1 : ex. Commerciaux]} & \textit{[N pers.]} & \textit{[Nouvelle interface, nouveaux process]} & \textit{[Faible / Moyen / Élevé]} \\
\rowcolor{cnIce}
\textit{[Équipe 2]} & \textit{[N pers.]} & \textit{[Description impact]} & \textit{[Niveau]} \\
\textit{[Équipe 3]} & \textit{[N pers.]} & \textit{[Description impact]} & \textit{[Niveau]} \\
\bottomrule
\end{tabularx}

\section*{2 — Réseau de Key Users}

\renewcommand{\arraystretch}{1.3}
\begin{tabularx}{\textwidth}{L{4cm} L{3cm} Y}
\toprule
\rowcolor{cnNavy}\th{Key User} & \th{Service} & \th{Rôle dans le déploiement} \\
\midrule
\textit{[Prénom NOM]} & \textit{[Service]} & Référent terrain, relais formation, remontée anomalies \\
\rowcolor{cnIce}
\textit{[Prénom NOM]} & \textit{[Service]} & \textit{[Rôle]} \\
\textit{[Prénom NOM]} & \textit{[Service]} & \textit{[Rôle]} \\
\bottomrule
\end{tabularx}

\section*{3 — Plan d'actions}

\renewcommand{\arraystretch}{1.3}
\begin{tabularx}{\textwidth}{C{1cm} L{4cm} Y C{2.5cm} C{2.5cm}}
\toprule
\rowcolor{cnNavy}\th{\#} & \th{Action} & \th{Description} & \th{Responsable} & \th{Timing} \\
\midrule
1 & Communication initiale & Annonce du projet, enjeux, planning & CP + DG & \textit{[Mois M+1]} \\
\rowcolor{cnIce}
2 & Ateliers de co-construction & Impliquer les Key Users dans les specs & CP + MOA & \textit{[Mois M+2]} \\
3 & Formations Key Users & Formation approfondie avant UAT & CP + ESN & \textit{[Mois M+4]} \\
\rowcolor{cnIce}
4 & Formation utilisateurs finaux & Sessions courtes, par groupe & Key Users & \textit{[Mois M+5]} \\
5 & Support post-déploiement & FAQ, hotline, permanence & DSI + Key Users & \textit{[Mois M+6]} \\
\rowcolor{cnIce}
6 & Mesure de l'adoption & Enquête satisfaction 30j post-go-live & CP + PMO & \textit{[Mois M+7]} \\
\bottomrule
\end{tabularx}

\section*{4 — Indicateurs d'adoption}

\renewcommand{\arraystretch}{1.3}
\begin{tabularx}{\textwidth}{L{5cm} C{2.5cm} Y}
\toprule
\rowcolor{cnNavy}\th{Indicateur} & \th{Cible} & \th{Modalité de mesure} \\
\midrule
Taux de participation aux formations & $\geq$ 90\% & Feuilles d'émargement \\
\rowcolor{cnIce}
Taux de satisfaction formation (enquête) & $\geq$ 80\% & Questionnaire post-formation \\
Taux d'utilisation de l'outil à J+30 & $\geq$ 70\% & Logs applicatifs \\
\rowcolor{cnIce}
Nombre d'incidents critiques J+30 & $\leq$ 5 & Log anomalies \\
Satisfaction utilisateur globale & $\geq$ 7/10 & Enquête NPS interne \\
\bottomrule
\end{tabularx}

\end{document}
